\begin{table}[t]
\centering
\caption{Maximum-likelihood parameter estimates for MS-TCM fit to the FRFR category condition. Values are MLE $[\text{95\% bootstrap CI}]$ (percentile method, $n_{\text{bootstraps}}=20$). Derived parameters are not fit independently (e.g.\ $w_S = 1 - w_G$).}
\label{tab:params}
\small
\begin{tabular}{@{}l l l p{0.52\linewidth}@{}}
\hline
Parameter & Symbol & Value $[\text{95\% CI}]$ & Description \\
\hline
\texttt{beta\_global} & $\beta_G$ & 0.950 \,[0.001,\,1.000] & Drift rate of the global temporal context. Higher values mean the global context updates more strongly toward each new event (and forgets older ones more quickly). \\
\texttt{beta\_storyline} & $\beta_S$ & 0.533 \,[0.241,\,0.551] & Drift rate of the storyline-specific context. A value near 0 means each storyline context behaves as a nearly static signature of the category; a value near 1 means within-storyline events overwrite one another as they are encoded. \\
\texttt{w\_global} & $w_G$ & 0.000 \,[0.000,\,0.543] & Weight on the global context in the composite encoding context. Higher values make the model more sensitive to global temporal contiguity (standard TCM sets $w_G = 1$). \\
\texttt{w\_storyline} & $w_S$ & 1.000 (derived) & Weight on the active storyline context in the composite. Higher values make category or storyline identity a stronger cue for recall. \\
\texttt{w\_global\_ret} & $w^{\mathrm{ret}}_G$ & 0.000 (derived) & Global-context weight at retrieval. Can differ from $w_G$ at encoding, reflecting task-driven reweighting of cues. \\
\texttt{w\_storyline\_ret} & $w^{\mathrm{ret}}_S$ & 1.000 (derived) & Storyline-context weight at retrieval. Increasing it models an instruction to recall within the same storyline. \\
\hline
\multicolumn{4}{l}{$\log L = -14067.2$, AIC $= 28140.4$, BIC $= 28160.1$; $n_{\text{participants}}=30$, $n_{\text{lists}}=480$, $n_{\text{recalls}}=5205$} \\
\end{tabular}
\end{table}
